% ============================================================
\section{Conclusions}
\label{sec:conclusions}
% ============================================================

This report has presented the design, implementation, and evaluation of
\textit{Ad LLM as a Service}, a working platform that delivers
sponsored product recommendations natively inside the responses of a
conversational AI. The platform was motivated by a structural mismatch
between the inventory of advertising formats inherited from the visual
web and the linear, multi-turn nature of chat-based interfaces, and
was scoped to demonstrate that this mismatch can be resolved without
either appending foreign ad blocks to the conversation or retreating
to keyword-style auctions over the most recent user message. The
sections below summarise what was built, partition the work between
the two workstreams, state the limitations of the current
implementation, and sketch the natural successors to the system as
delivered.

% ------------------------------------------------------------
\subsection{Summary of Contributions}
\label{subsec:summary-contributions}
% ------------------------------------------------------------

The platform contributes a set of components that are, individually,
adaptations of known techniques, but whose joint realisation inside a
single end-to-end system has not, to the best of our knowledge, been
published before. We summarise the contributions as follows.

\begin{enumerate}[label=(\roman*), leftmargin=2.5em]
  \item \textbf{An end-to-end architecture for LLM-native
    conversational advertising, spanning both text and image
    modalities.} A single FastAPI orchestrator serves a chat pipeline
    that produces ad-augmented textual replies and an image-generation
    pipeline that produces ad-augmented imagery, both reusing the
    same context extraction and matching infrastructure.

  \item \textbf{A specialised RAG pipeline for ad matching.} Semantic
    retrieval over a \texttt{pgvector} index is combined with a
    business-rule re-ranker that enforces campaign eligibility, daily
    and total budget caps, per-session frequency capping, and a
    weighted composite score in which semantic similarity, bid price,
    and purchase signal are interpretable and auditable.

  \item \textbf{Prompt-engineered natural ad integration with
    mandatory disclosure.} Ad injection is realised as
    prompt-level conditioning of the response-synthesis model rather
    than as post-hoc text rewriting, and every commercial mention is
    accompanied by the literal \texttt{[Sponsored]} marker required by
    Federal Trade Commission guidance.

  \item \textbf{Multi-dimensional engagement measurement via an
    in-the-loop LLM judge.} A dedicated judge call returns, per
    follow-up turn, a six-field engagement vector comprising an
    engagement-type label, a continuous engagement score, a
    naturalness score, an estimated purchase-proximity, a sentiment
    polarity, and a topic-match flag, which jointly form a
    higher-fidelity signal than any single click-or-not-click
    measurement.

  \item \textbf{Closed-loop adaptive optimisation of per-ad
    presentation context.} Each advertisement carries a versioned
    presentation-guidance document that is rewritten by an optimiser
    LLM in response to aggregated engagement metrics, under
    minimum-engagement and minimum-interval guards, with a
    natural-language rationale and a metrics snapshot persisted for
    every revision.

  \item \textbf{Organic product placement in generated imagery.} The
    image pipeline rewrites the user's prompt to place the matched
    product diegetically into the requested scene and, when the
    advertisement carries a reference photograph, dispatches the
    rewritten prompt to the image-edit endpoint so that the rendered
    product is visually grounded in the real article rather than
    reconstructed from a textual prior.

  \item \textbf{A full advertiser portal with campaign, budget, and
    ad management.} The portal exposes CRUD operations over
    advertisers, campaigns, and ads under a strict URL hierarchy,
    enforces a \texttt{draft}/\texttt{active}/\texttt{paused}/\texttt{completed}
    campaign lifecycle with auto-completion on budget exhaustion, and
    accelerates ad creation via a URL-based extraction flow that
    scrapes a product page and pre-fills the ad fields as an editable
    draft.

  \item \textbf{Ethical safeguards realised as structural invariants
    rather than soft filters.} Sensitivity-aware gating at the
    context-extraction stage short-circuits the matching pipeline
    before any retrieval or budget accounting work is performed, and
    this decision cannot be overridden by bid, similarity, or budget
    on any downstream branch.
\end{enumerate}

% ------------------------------------------------------------
\subsection{Workstream Contributions}
\label{subsec:workstream-contributions}
% ------------------------------------------------------------

The project was delivered by us developing in
parallel against the frozen interface contract of
Section~\ref{subsec:interface-contracts}. The \textbf{backend and
infrastructure} workstream owned the FastAPI orchestrator, the
session manager, the PostgreSQL schema (including the \texttt{ads},
\texttt{campaigns}, \texttt{daily\_spend\_log}, \texttt{events}, and
\texttt{ad\_context\_history} tables), the ad matching pipeline with
its budget-gated eligibility check, the click-tracking and
spend-accounting transaction, the platform and advertiser analytics
APIs, the React dashboards, the advertiser portal including the
campaign lifecycle endpoints, the image-generation route, and the
frontend integration for both the chat and image surfaces. The
\textbf{AI and LLM} workstream owned the context extractor and its
Pydantic schema, the embedding pipeline on both the ingestion and
query sides, the response-synthesis prompt and its context-guide
injection, the engagement judge, the initial-context generator invoked
at ad ingestion time, the periodic context optimiser, the placement
prompt rewriter used in the image pipeline, the URL-based product
extraction flow, and the LLM-as-judge offline evaluator used in the
quality-assurance framework. Each workstream was able to exercise its
own invariants against a typed stub of the other, so that system
integration was a short merge step rather than a protracted debugging
phase.

% ------------------------------------------------------------
\subsection{Limitations}
\label{subsec:limitations}
% ------------------------------------------------------------

The platform as delivered is a minimum viable product and carries a
number of limitations that are worth stating explicitly. The chat and
image pipelines depend on third-party LLM and image-generation APIs,
which subjects the serving path to the latency, cost, and availability
envelope of those providers rather than to a self-hosted baseline.
The re-ranking weights $(0.5, 0.3, 0.2)$ are hand-tuned
hyperparameters rather than coefficients learned from click or
engagement data, and although they are explicable at the level of the
design objective, they are not optimal in any formal sense. The
engagement vector on which the feedback loop operates is produced by
a single LLM judge call per follow-up turn and is therefore subject
to drift as the underlying judge model is revised by its provider;
we do not currently run a second judge in parallel or calibrate the
score against held-out human labels. The adaptive context optimiser
treats each advertisement as a single target and produces one guide
per ad, which precludes segment-conditional presentation strategies
in which the recommended register or level of specificity depends on
the user's inferred context. The image-generation surface lacks any
automated visual-disclosure mechanism beyond the ad-metadata payload
returned alongside the rendered image, and the trademark and
trade-dress implications of rendering a reference-grounded product
into a user-specified scene have not been formally assessed. The ad
inventory used during development was small, so claims about
inventory-scale behaviour beyond the order-of-magnitude targets of
Section~\ref{subsubsec:ad-inventory-store} rest on analytical rather
than empirical evidence. Finally, the platform has not been evaluated
with real advertiser traffic or under an A/B testing protocol, and
the engagement-rate claims quoted elsewhere in the report are based
on the internal regression set rather than on a production user base.

% ------------------------------------------------------------
\subsection{Future Work}
\label{subsec:future-work}
% ------------------------------------------------------------

The limitations above map directly onto a roadmap of successor
problems. We group them by the component they extend.

\begin{itemize}[leftmargin=1.8em]
  \item \textbf{Attention and learned ranking.} Extend the engagement
    judge with a dwell-time and read-through estimator, and train a
    predicted-CTR model on the accumulated click and engagement log to
    replace the hand-tuned re-ranking weights with learned
    coefficients under appropriate exploration.

  \item \textbf{Segment-conditional ad context.} Support multiple
    presentation guides per advertisement, selected at synthesis time
    by the extracted context object, so that the register and
    specificity of a sponsored mention can adapt to the user's
    inferred state without requiring a separate creative.

  \item \textbf{Judge ensembling and calibration.} Run a second,
    architecturally distinct LLM judge in parallel on a sampled
    subset of turns, periodically reconcile the two scores against a
    small human-labelled set, and treat the agreement rate itself as
    a monitored metric.

  \item \textbf{Visual disclosure for generated imagery.} Introduce a
    watermark, caption-level tagging, or provenance-metadata channel
    on images that carry a sponsored product, so that the visual
    surface acquires the analogue of the \texttt{[Sponsored]} marker
    that the text surface already enforces.

  \item \textbf{Principled evaluation under within-session
    dependency.} Replace the current regression-set methodology with
    a bandit-style evaluation that respects the observation of
    Section~\ref{subsec:qa-framework} that serving decisions on one
    turn shape context on subsequent turns within the same session,
    and thereby violate the independence assumptions that
    conventional A/B tests rely on.

  \item \textbf{Tone adaptation.} Make the register of the sponsored
    mention responsive to the surrounding conversational register, so
    that a casual chat thread does not suddenly receive a
    brochure-style recommendation and a technical thread does not
    receive a colloquial one.

  \item \textbf{Brand safety and content filtering.} Add an explicit
    topic classifier and blocklist upstream of response synthesis,
    complementing the existing sensitivity gate with an
    advertiser-configurable brand-safety layer.

  \item \textbf{Commercial integration.} Connect the campaign spend
    ledger to a billing and invoicing subsystem so that the
    platform's accounting of charged impressions becomes the
    authoritative input to advertiser billing rather than an internal
    diagnostic.

  \item \textbf{Domain-specific context extraction.} Fine-tune the
    context extractor on domain-specific dialogue data so that the
    LLM call can be replaced with a smaller, cheaper, and
    more targeted model without surrendering schema fidelity or
    receptivity-gate sensitivity.
\end{itemize}

% ------------------------------------------------------------
\subsection{Reflections}
\label{subsec:reflections}
% ------------------------------------------------------------

Three observations from the build process are worth recording
explicitly, not as project artefacts but as notes for future work of
a similar shape.

The first concerns the value of a frozen interface contract at the
start of a two-workstream project. The four function signatures and
four dataclasses of Section~\ref{subsec:interface-contracts} were
agreed before any production code existed, and every subsequent
design decision that might otherwise have crossed the seam between
workstreams was instead resolved inside one of the two boundaries.
The cost of this discipline was a single unrushed session in Week~1;
the benefit was that system integration never became a blocking
activity.

The second concerns prompt engineering. The response-synthesis prompt,
the engagement-judge prompt, and the context-optimiser prompt each
passed through a double-digit number of revisions during development,
and the revisions that produced the largest behavioural change were,
without exception, those that tightened a constraint rather than
those that added new instructions. The prompt that yielded the most
natural ad integration was the one that specified what the model was
permitted to omit, not the one that specified what it was required to
say.

The third concerns the tension between user experience and
monetisation that animates the entire problem. It is tempting, under
deadline pressure, to treat this tension as a slider and to argue
about where to set it. The present project took the opposite stance:
the helpfulness of the response was treated as a precondition that
every ad-related component had to respect, not as a quantity that
could be traded off against commercial yield. Every ethical
safeguard, from the sensitivity gate in the context extractor to the
final-check pattern before response synthesis, is a structural
encoding of this stance, and the fact that each could be implemented
without sacrificing the interactive latency budget suggests that the
tension was less fundamental than it first appeared.

Taken together, the contributions, limitations, and successor
problems above define a platform that is, we believe, a faithful
instantiation of the thesis with which this report began: that
advertising in conversational AI is not a porting problem but a
format problem, and that the format it calls for is one in which the
advertisement is written by the same model that writes the answer,
within the same conversation, under the same constraints of
helpfulness, disclosure, and care.